\documentclass[10pt]{article}
\usepackage[cmex10]{amsmath}
\usepackage{amsthm}
\DeclareFontShape{OMX}{cmex}{m}{n}{<->cmex10}{}
\DeclareMathSizes{10}{10}{7}{7}
\DeclareMathSizes{9}{10}{7}{7}
\setlength{\textwidth}{6in}
\setlength{\oddsidemargin}{0.25in}
\setlength{\evensidemargin}{0.25in}
\setlength{\textheight}{8.85in}
\setlength{\topmargin}{-0.45in}
\newtheorem{theorem}{Theorem}
\newtheorem{lemma}{Lemma}
\newtheorem{corollary}{Corollary}
\newtheorem{proposition}{Proposition}
\title{A 24-edge lower bound for rank-six intersecting Ryser counterexamples}
\author{\vspace{-1.5\baselineskip}}
\date{November 26, 2026}

\begin{document}
\maketitle

\begin{abstract}
Every six-partite, six-uniform intersecting hypergraph with cover number six
has at least 24 edges. This improves the reported lower bound of 22.
Deleting two vertex stars leaves at least nine edges, and deleting one
leaves at least 14. Integer degree-profile enumeration and a side congruence
exclude 22 edges. Transportation and overlap constraints reduce the
23-edge case to three systems of overlap masks. Exact integer inequalities
exclude all three. At 24 edges the maximum degree is at most nine.
We also exclude counterexamples with order-six rotation symmetry
through 43 edges and extensions of at least twelve unchanged rows
of the affine plane of order five.
\end{abstract}

\begin{theorem}
If $H$ is a finite, simple, intersecting six-partite, six-uniform hypergraph
with covering number $\tau(H)=6$, then $|E(H)|\geq24$.
\end{theorem}

\section{Deleting stars}

Let $m=|E(H)|$ and write $E(v)$ for the edges through $v$. Ignore isolated
vertices throughout. Each part is a cover and has at least six vertices.
No edge has a degree-one vertex, since removing that vertex from the edge
would leave a five-vertex cover. More generally, for distinct vertices
$u,v$ one cannot have $E(u)\subseteq E(v)$. Any edge through $u$ would
contain $v$ and could omit $u$ from its cover.

Deleting all edges through a set of $k$ vertices leaves an intersecting
hypergraph of cover number at least $6-k$. Its edges can be covered in
pairs, choosing one shared vertex for each pair. Hence deleting three,
four, or five stars leaves at least five, three, or one edges, respectively.

\begin{lemma}
Deleting any two distinct stars leaves at least nine edges.
\end{lemma}

\begin{proof}
Write $J$ for the residual hypergraph, with $\tau(J)\geq4$.
The pair cover shows $|E(J)|\geq7$. At seven edges a vertex of degree at
least three would leave at most four other edges and give a three-vertex
cover. If every degree is at most two, each edge needs all six of its
vertices to have degree two so it can meet its six peers. This would
partition the seven edges into pairs on every side, which is impossible.

At eight edges, $\tau(J)=4$, and Lemma~2.1 of Abu-Khazneh and Pokrovskiy
\cite{akp} gives a degree-three vertex on every side and a pair of edges
sharing two degree-three vertices. If the two selected stars are in one
side of $H$, that side retains at least four disjoint blocks, each of
unchanged degree at least two. With exactly eight edges these four blocks
all have degree two, which contradicts the lemma.

If the stars are on different sides, each of their two sides retains at
least five vertices. Indeed the other vertices on either side are
disjoint from the deleted star in that side and their stars cannot be
contained in the other deleted star. In $J$, the largest degree is at most
three, since a degree-four vertex and two pair covers would give a
three-vertex cover. Two degree-three vertices on one side would also give
such a cover. Thus each side contributes at most five pair intersections,
and a side with at least five vertices contributes at most four. The total
is at most $2\cdot4+4\cdot5=28$. Since the 28 edge pairs must intersect
and one pair intersects twice, this is impossible.
\end{proof}

\begin{lemma}
Deleting any one star leaves at least 14 edges.
\end{lemma}

\begin{proof}
Denote the residual by $J$. Its cover number is at least five, so it has
at least 13 edges by Aharoni, Bar\'at, and Wanless \cite[Theorem~2.7]{abw}.
Suppose it has exactly 13.
Every vertex outside the side of the deleted star survives, by the
noncontainment of stars. In the affected side, at least five vertices
survive. The other five sides each have at least six. No edge of $J$
has two degree-one vertices, since removing both would give a cover
of size four. Hence at most 13 vertices of $J$ have degree one.

In each side, the sums of the largest $k$ degrees are bounded by
$4,8,10,12$ for $k=1,2,3,4$. Deleting a further $k$ stars in that
side is equivalent to deleting $k+1$ original stars. Apply the first
lemma for $k=1$ and the pair cover bounds for $k=2,3,4$.
The integer partitions of 13 into at least five positive degrees
that meet these prefix bounds give 35 side profiles. Enumerating
six-profile multisets with at most one side of size five and at most
13 degree-one vertices gives zero systems whose pair score
$\sum_v\binom{d(v)}2$ is at least $\binom{13}2=78$.
The supplied verifier checks this finite enumeration.
\end{proof}

\section{Pair intersections at 22 edges}

Let $p_i$ be the nonincreasing degree profile of side $i$ of $H$.
Its entries are at least two, its length is at least six, and the sums of
its first one through five entries are at most
\[
m-14,\quad m-9,\quad m-5,\quad m-3,\quad m-1.
\]
Write $s(p)=\sum_{a\in p}\binom a2$. Every two edges meet, so
$R=\sum_i s(p_i)-\binom m2$ is nonnegative.

Every edge contains at most one degree-two vertex. If an edge contained
two, dropping both from its cover would uncover at most their two
other incident edges. Intersectingness covers these with one vertex.
Thus the sum of the numbers of degree-two entries of the six profiles
is at most $\lfloor m/2\rfloor$.

An edge pair cannot agree in five sides. An edge avoiding the five
common vertices would have to meet the pair in two distinct vertices
from the same side. Consequently the five common vertices would
already form a cover.
Let $n_k$ count pairs meeting in exactly $k$ sides, for $1\leq k\leq4$.
For distinct sides $i,j$, let $T_{ij}$ count edge pairs agreeing in both
sides, and put $T=\sum_{i<j}T_{ij}$. Then
\[
 R=n_2+2n_3+3n_4,\qquad T=n_2+3n_3+6n_4,\qquad 2R-T=n_2+n_3.
\]

For profiles $a,b$ of two sides, let $c_{uv}$ be the number of
edges through the vertices of degrees $a_u,b_v$. Its matrix has
row sums $a$ and column sums $b$. The two-star lemma and
noncontainment give
\[
 \max(0,a_u+b_v-(m-9))\ \leq c_{uv}\ \leq\min(a_u,b_v)-1 .
\]
Minimize $\sum_v\binom{c_{uv}}2$ for each row under these bounds and
its row sum. The sum of the minima bounds $T_{ij}$ from below.
Repeat with columns and take the larger bound $L_{ij}$. This uses
only one-dimensional integer enumeration. A profile system can
survive only if $\sum_{i<j}L_{ij}\leq2R$.

The verifier enumerates all side profiles for each $m$ from 13
through 22 and all six-profile multisets subject to these conditions.
The lower endpoint follows from $\tau(H)\geq5$ and \cite{abw}.
No system survives for $m\leq21$. At $m=22$ there are 64 individual
profiles and five surviving systems. All five have exactly eleven
degree-two vertices, so every edge contains one.

In this case a degree-three vertex cannot be used in a multiple
intersection. For suppose $e$ meets another edge at its degree-three
vertex $v$ and at a second vertex. Drop $v$ and the degree-two vertex
of $e$ from the cover $e$. The edge sharing two vertices either
stays covered or contains both deleted vertices. In either case
at most two edges remain uncovered. One more vertex covers them.
The contradiction implies $c_{uv}\leq1$ whenever either degree is three.
Recomputing the row and column bounds eliminates two systems.

\begin{center}
\begin{tabular}{c|c|c|c}
$R$ & $\sum L_{ij}$ & $2R-\sum L_{ij}$ & One side lower bound\\
\hline
20 & 40 & 0 & 10\\
18 & 35 & 1 & 13\\
15 & 30 & 0 & 8
\end{tabular}
\end{center}

If $2R-\sum L_{ij}=0$, then $n_2+n_3=0$. This excludes $R=20$,
which is not divisible by three. If $2R-\sum L_{ij}=1$ and $R$ is
divisible by three, one again has $n_2=n_3=0$. Indeed the only
other possibilities give remainder one or two in $R$.
When all multiple intersections have size four, each side has
$\sum_{j\ne i}T_{ij}$ divisible by three. Each size-four
intersection using that side contributes three to this sum.
For $R=18$ the indicated side sum lies between 13 and 14.
For $R=15$ it equals eight. This excludes 22 edges.

\section{The 23-edge case}

Suppose that $m=23$. Let $t$ be the number of degree-two vertices
and keep the notation $R$, $T_{ij}$, and $T$ from the previous section,
now with $\binom{23}{2}=253$. Every degree-two star contains two
edges meeting only at its center. Indeed, if they also met elsewhere,
deleting its center from one of the two edges would leave a five-cover.
Consequently a codegree involving a degree-two vertex is at most one,
and is zero between two distinct degree-two vertices. If $t=11$, every
pair has at least one endpoint with a degree-two vertex, so a codegree
involving degree three is also at most one, by the earlier argument.

\subsection{Profile and star reductions}

Enumerate the six profiles with at least six positive entries, all
between two and nine, and with the sum of the two largest at most 14.
The score must be at least 253 and $t\leq11$. For every pair
of profiles, minimize the collision score across the entire codegree
matrix with given margins and the established cell bounds. After
subtracting the entrywise lower bounds, encode each allowed unit into
a cell by a parallel arc of unit capacity. The successive costs
for adding to a cell of current size $a$ are $a,a+1,\ldots$.
This makes the minimum flow cost plus the initial cost the exact
minimum of $\sum_{u,v}\binom{c_{uv}}2$. Reject if the sum of the
15 minima exceeds $2R$. Repeat with the degree-three cap when $t=11$.

Next consider a particular vertex $v$ of degree $d$. Its codegree
row into another side with profile $q$ is a vector $a$ of sum $d$,
within the same cell bounds. The residual degrees on that side are
$q-a$. The sums of the largest one through four residual
degrees on any side are at most $23-d-9$, $23-d-5$,
$23-d-3$, and $23-d-1$ by the star deletion bounds. The side of $v$
is treated by deleting the entry $d$ from its profile. The residual
cover number is at least five, so at most $23-d$ residual vertices
have degree one. Moreover a row through $v$ uses five other vertices
that cover the residual. At most one of them has residual degree one,
since two such vertices can be replaced by a single vertex covering
their at most two residual edges. Thus the total incidences of the
original star with vertices of residual degree one are at most $d$.
The residual score is at least $\binom{23-d}{2}$, and the cross score
$\sum_{w}c(v,w)(d(w)-c(v,w))$ is at least $d(23-d)$, since every
edge through $v$ meets every edge outside that star.

Enumerate all allowable vectors $a$ independently for each other side.
For each tuple of degree-one count, degree-one incidence in the star,
and residual score, it suffices to retain the largest cross score.
Join these finite tables over the sides. Apply the test to a largest
vertex on every side whose maximum is at least eight. In a second
pass use a threshold of seven and also track collision surplus.
For each candidate $a$ in a pair matrix, remove its fixed row and
minimize the remaining collision score separately by rows and by
columns with the resulting margins and cell bounds. Add the fixed
row cost, take the larger lower bound, and charge any excess above
the unrestricted minimum against $2R$ minus the sum of the pair minima.
Only the five pairs touching the side of $v$ accrue these charges,
so no surplus is counted twice for a test of one vertex.

\subsection{Pair bounds with subset unions and overlap patterns}

Strengthen a pair matrix by applying the bounds on the union of
two through five vertex stars to every choice of vertices from
its two sides. The sum of the selected degrees minus the sum of
entries in the selected rows and columns is the size of the union.
The parallel unit-arc formulation, with these extra linear
constraints, gives a rational relaxation. Every lower bound used
for pruning is validated by an exact rational dual inequality
and rounded up, since the collision count is an integer.

An edge pair meeting in a set $I$ of two, three, or four sides costs
$|I|-1$ from the redundancy budget and contributes one to each required
pair of side indices in $\binom{I}{2}$. A finite recursion chooses
patterns to meet the lower bounds on the 15 side pairs at cost at most
$R$. It branches on a side pair with positive remaining demand and
tries every pattern containing it. When $t=11$, a pattern of size four
must omit at least one side with a degree-two vertex, since at least
one endpoint carries such a vertex and its degree-two side cannot be
among the overlaps.

\begin{center}
\begin{tabular}{l|r}
Necessary conditions applied & Remaining profile systems\\
\hline
Total score and degree-two budget & $1{,}756{,}829$\\
Exact pair transportation & $9{,}614$\\
Degree-three cap when $t=11$ & $5{,}527$\\
First star test & $537$\\
Star test with collision surplus & $126$\\
Subset-union transportation bounds & $72$\\
Overlap pattern budget & $17$
\end{tabular}
\end{center}

\subsection{Row bounds and the last profile}

Every vertex of an edge $e$ has a private neighbor edge, meaning
an edge sharing just that vertex with $e$. Otherwise the vertex could
be omitted from $e$ to obtain a five-cover. If $e$ has a degree-two
vertex, each of its other vertices needs at least two private
neighbors. If such a vertex had at most one, deleting it and the
degree-two vertex would leave at most two uncovered edges, since the
unique neighbor at the degree-two vertex is already private.

For a prospective row of degrees $a_1,\ldots,a_6$, its total excess
from repeated intersections is $\sum_i a_i-28$. Suppose the row has
$b_k$ neighbors with $k$ common vertices. The excess is
$b_2+2b_3+3b_4$. In a side of degree $a_i$, at most $a_i-1$ neighbors
occur there, minus the required number of private neighbors.
Call the resulting capacity for repeated neighbors $u_i$, setting
it to zero at degree two and wherever the degree-three restriction
applies. The $b_4$ neighbors need $4b_4$ incidences in total and
can use at most $\min(u_i,b_4)$ in side $i$. Apply the same inequality
to the combined groups of fourfold and threefold neighbors and to all
repeated neighbors. Enumerating $b_2,b_3,b_4$ gives a lower bound on
$b_2+b_3$ for each allowed degree tuple. Summing over the 23 rows,
with the exact margins for each degree in each side, must give
at most $2(2R-\sum_{i<j}L_{ij})$. The factors of two arise because
each pair has two endpoints. Verified rational linear inequalities
exclude 16 of the 17 profile systems.

The remaining system, in order, is
\[
\begin{gathered}
(8,6,3,2,2,2),\quad (7,7,3,2,2,2),\quad
(8,5,3,3,2,2),\\
(8,5,3,3,2,2),\quad
(8,4,3,3,3,2),\quad (8,3,3,3,3,3).
\end{gathered}
\]
Here $t=11$, $R=23$, and the sum of the strengthened pair minima
is 44. Thus there are seven fourfold pairs and either one threefold
pair or two twofold pairs. Enumeration of the side masks gives no
systems with two twofold pairs and three systems with one threefold pair.

Consider the graph whose vertices are the edges of $H$ and whose
adjacencies are fourfold intersections. If two adjacent edges each
have a degree-two vertex, it must be in the same side. To see this,
remove their four common vertices. An avoiding edge intersects both
in one of two cross combinations on the remaining sides. Both
combinations must occur, since otherwise a fifth vertex suffices
to cover. If the degree-two vertices were on different sides, one
cross combination would contain both, which is impossible.

The fourfold graph has no triangle. A putative triangle has at
least two edges of $H$ with degree-two vertices, in the same side
$i$ by the preceding argument. Any third edge shares neither of
these degree-two labels, since its pair would otherwise meet only
there. All three pairs therefore disagree in side $i$. Their
projections to the other five sides have pairwise Hamming distance
one, so all three projections differ in the same side $j$.
For a triangle edge carrying a degree-two vertex, its unique mate
would have to supply two different $j$-values in the two cross
combinations with the other triangle edges. A single edge cannot do so.

Two neighbors of a row in this triangle-free graph have at least
two agreements, because each differs from the row in two sides.
Their pair therefore must be the unique threefold pair. No row
can have three fourfold neighbors, and at most one row can have
two. If two distinct rows were common to the endpoints of the
threefold pair, those two rows would supply a second pair meeting
in two or three sides. The row with two fourfold neighbors is
outside the unique threefold pair. The two fourfold masks through
that row intersect in precisely the threefold mask, since the two
sets of differing sides must have exactly one common side.

For each of the three overlap-mask systems, enumerate a row's
degree tuple and the repeated-pair labels incident to that row.
Apply the preceding restrictions. For every proper nonempty set
$S$ of side indices, deleting the row's vertices in the other
sides leaves at least $1,3,5,9,14$ edges when $|S|=1,\ldots,5$.
These remaining edges are exactly the row's private neighbors in
$S$ and its repeated neighbors whose entire masks lie in $S$.
This produces $78,71,78$ allowable local rows respectively.

Each mask label needs two endpoints. The degree counts on each
side are fixed. For each of the three systems the verifier checks
an integer linear functional that is nonpositive on every
allowable local row and takes value respectively 55, 53, and 55
on the required aggregate margins. The aggregate therefore cannot
be realized. This excludes 23 edges.

\section{Restrictions at 24 edges}

\begin{proposition}
If a rank-six counterexample has 24 edges, then every vertex has
degree at most nine.
\end{proposition}

\begin{proof}
The one-star bound gives maximum degree at most ten. Suppose
that $v$ has degree ten and delete the ten edges through it, giving
a residual hypergraph $J$ with 14 edges and cover number at least five.
Its largest degree is at most five by the two-star bound. The sums
of the three and four largest degrees in a side are at most 11 and
13 by the four- and five-star bounds. On the side of $v$, at least
five nonisolated vertices remain and their degrees are unchanged.
All other sides retain at least six vertices by noncontainment.

There are at most ten degree-one vertices of $J$. To see this, the
other vertices of each deleted edge form a five-cover of $J$. At
most one vertex in that cover can have residual degree one. If
there were two, a single vertex could cover their at most two edges
and reduce the cover to four. Every residual degree-one vertex
occurs on at least one deleted edge because its degree in the
original hypergraph was at least two. Counting these incidences gives
the claimed bound.

Let $k$ be the number of degree-five vertices, $q\in\{0,1\}$
indicate whether one is on the side of $v$, and $k'=k-q$.
No side contains two degree-five vertices, because deleting their
stars in $J$ leaves at least five edges. Let $X_4$ count the
degree-four vertices, $W$ the degree-one vertices, and $K$ the sum
of the six side sizes in $J$. For $1\le d\le5$,
\[
 \binom d2-[d=1]=2d-3+[d=4]+3[d=5].
\]
The score $S$ that counts pair intersections with multiplicity
therefore satisfies
\[
 S=168-3K+W+X_4+3k\ \leq\ 73+X_4+3k,
\]
since $K\ge35$ and $W\le10$.
The side of $v$ has at most two degree-four vertices if $q=0$
and none if $q=1$. Indeed three such vertices violate the bound on
the three largest degrees, and a degree five and a degree four
would leave fewer than six degrees for the other three vertices
on that side. Each other side has at least six vertices. A side
there with a degree-five vertex has at least one singleton by
degree count, and a degree-four vertex on that side requires at
least two more singletons. A side there without a degree-five
vertex can have a first degree-four vertex freely, while a second
requires at least two singletons. Thus
\[
 X_4\le (5-k')+2(1-q)+\left\lfloor\frac{10-k'}2\right\rfloor.
\]
The resulting upper bounds on $S$ for $k=0,\ldots,6$ are
$85,86,88,89,91,92,93$. Intersectingness requires $S\ge91$.

Every two degree-five stars in $J$ intersect, since their union
has size at most nine. No three can have all their pairwise
codegrees equal to one. Their union would contain at least twelve
edges, giving a four-cover after covering the remaining two edges
with one vertex. Put an edge between two degree-five vertices if
their codegree is at least two. Its complement is triangle-free
and has at most $\lfloor k^2/4\rfloor$ edges by Mantel's bound.
Hence these stars alone contribute at least
$\binom k2-\lfloor k^2/4\rfloor$ to the two-side collision count $T$.
Each pair of residual edges still meets in at most four vertices,
so $T\le2(S-91)$. For $k=4,5,6$ the respective lower bounds on
$T$ are $2,4,6$ and the upper bounds are $0,2,4$.
This is a contradiction in every case.
\end{proof}

Two further constraints follow from the private pairs at degree-two
vertices. In an $m$-edge counterexample with $t$ degree-two vertices,
the resulting $t$ pairs are disjoint and leave $B=m-2t$ unpaired
edges. If part $i$ contains $t_i$ of those vertices, any other vertex
in that part is in at most one edge of each pair and in none of the
$t_i$ pairs from its own part. Thus its degree is at most $t-t_i+B$.

For distinct vertices $v,w$ of degrees at least three, suppose that
$d(w)-c(v,w)=1$. Every edge containing both $v,w$ is unpaired:
an original degree-two vertex on such an edge would have residual
degree one after deleting the star of $v$, since its private mate
avoids $v$. Together with $w$ this would give two residual
singletons in a five-cover. Therefore $c(v,w)\le d(w)-2$ whenever
$d(w)-1>B$, and the corresponding inequality holds when interchanging
$v,w$. At 24 edges, applying the mate and codegree constraints to
the degree and pair-transport enumeration leaves 265,367 degree
systems prior to the joint residual-star tests.

\begin{proposition}
Deleting any active vertex star from a rank-six counterexample leaves
two edges that meet in at least two vertices.
\end{proposition}

\begin{proof}
The residual has cover number at least five. The other five vertices
of any deleted edge cover it, so its cover number is exactly five.
Noncontainment of stars implies that every active vertex other than
the deleted vertex remains active. The residual therefore has at least
35 active vertices. If it were linear, Theorem 4.1(1) of \cite{fhmw}
would embed it into the truncated projective plane of order five,
which has only 30 vertices.
\end{proof}

In particular, form a graph on the edges of $H$ by joining pairs that
meet in at least two vertices. No vertex star of $H$ is a vertex cover
of this graph. Its vertex cover number is at least three: any one or
two edges of $H$ lie in a common star. Each one-star residual must
also have at least one excess intersection incidence over its number
of pairs.

\section{Rotation symmetry}

\begin{proposition}
Suppose an automorphism of order six cycles the six parts of $H$.
If $H$ has at most 43 edges or at most ten active vertices per part,
then $\tau(H)\le5$.
\end{proposition}

\begin{proof}
Transport the labels of one part to the others by powers of the
automorphism. Its sixth power is the identity, so the automorphism
becomes the cyclic rotation of row coordinates. When the common
part size is at most ten, pad the alphabet to ten and enumerate the
1,000,000 rows. In a counterexample distinct rows must agree in at
least one and at most four coordinates. The latter bound follows
from star noncontainment, or from the five shared vertices covering
every row. After this test within each orbit, 16,660 orbits remain.
Two distinct orbits can coexist only if all their cross pairs obey the
same test. To check compatibility it suffices to fix a representative
in the first orbit and test every row in the second.

A potential counterexample contains an orbit. A common relabeling of
all six parts commutes with rotation. The equality pattern of each
rotation of a representative therefore yields a canonical candidate
root by normalization in order of first occurrence. The program
starts once at each of the resulting 17 roots with all compatible
orbits available. To find a cover of size at most five, it branches
on the six vertices of a currently uncovered row with one fewer
available choice at each step. If the chosen orbits have a five-cover,
some remaining orbit in a counterexample must contain a row missed
by that cover. Order these orbits and branch on the first one that the
counterexample uses. Each branch retains compatible orbits that meet
the cover and compatible later orbits that break it. A five-cover of
the union of all available orbits excludes the entire branch.
The integer search in \texttt{edge24/cyclic.cpp} exhausts all roots.

Now suppose the common active part size $q$ is at least eleven and
$m\le43$. The main theorem gives $m\ge24$. Each part has the same
number of degree-two vertices, at most $b=\lfloor m/12\rfloor$ because
the degree-two pairs are disjoint. Every other active vertex has
degree at least three. This requires $m\ge33-b$, which excludes
$m<31$ in the stated range. For a fixed number of twos, convexity
maximizes the pair count by assigning all excess degree to one
vertex. The resulting upper bound increases on reducing the side
size to eleven and increasing the number of twos to $b$. The total
over all six sides is at most
\[
 6\left(\binom{m-30+b}{2}+30-2b\right).
\]
Subtract this from $\binom m2$. For each constant value of $b$ the
difference is a concave quadratic in $m$. At the endpoints 31, 35,
36, and 43 it equals 291, 313, 270, and 39, respectively. So the
pair score falls short of the number of pairs throughout the range.
\end{proof}

\section{Extensions of an affine-plane core}

Give each point of the affine plane over the field of five elements
the row of its six incident lines, one per direction. These labels are old and any
other labels are called new.

\begin{proposition}
An intersecting extension containing at least twelve unchanged
affine rows, in their original parts, has a five-cover.
\end{proposition}

\begin{proof}
It suffices to keep any twelve of the specified rows as the core.
A counterexample needs at least one new label in every direction.
Each row containing new labels must meet the core using its old
lines. Two nonparallel lines cover only nine points, so at least
three labels in the row are old. Enumerate all choices of one to three
omitted directions and all old line choices in the other directions.
For each union, mark every twelve-point subset of it as admitting new
labels in the omitted directions. There are 5,169,500 subsets that
admit new labels in all six directions. The group of affine linear
maps and translations partitions these into 482 orbits.

For each representative, use one symbolic new label per direction.
Actual equalities between new labels imply symbolic equalities. A
cover using only old labels therefore applies to every realization
of the covered symbolic types. A direction with fewer than two
singleton-new types and no multi-new type is impossible. A new
label would lie in at most one edge, since all five other labels are
fixed by the sole allowed type. This excludes 17 orbits. If a
direction has exactly two singleton-new types and no multi-new type,
both types must occur. Of the remaining orbits, 124 have an old-label
five-cover after filtering for compatibility with all such mandatory
types.

For the remaining 341 orbits, start with the core and the mandatory
types. Branch on the first available type breaking the current
old-label five-cover, retaining only compatible future types. The
search rejects a branch if all remaining types have an old-label
five-cover, if a needed direction can no longer support a new label
of degree at least two, or if a chosen old label can no longer occur
twice. For the last test an old-only type contributes one possible
occurrence, and a type with a new label contributes two, which
suffices for a threshold of two. The program never limits how many
actual rows can realize a type with new labels. Every actual extension
would therefore survive some branch. The exact search in
\texttt{edge24/affine-core.js} exhausts all branches in 18,645 nodes.
\end{proof}


\section{An affine-plane construction that cannot give a witness}

\begin{corollary}
Start with the affine plane of order five and retain any subset of its
points. In each of the six parallel directions, merge any collection of
whole line traces into blocks and partition the remaining individual
line traces into blocks. Use each retained point as an edge, with its
six block memberships giving its vertices. If the resulting hypergraph
is intersecting, its cover number is at most five.
\end{corollary}

\begin{proof}
If it has cover number six, at most one point was deleted by the
theorem. All five traces in each direction are nonempty.
No block can merge three lines because its deletion would leave at
most ten edges. Two merged blocks in one direction leave at most five
edges outside them, which can be covered by three more vertices.
Thus each direction has at most one merged block of two line traces.

Every block has size at least two. Splitting an intact five-point
line would therefore split it into sizes two and three. The six cross
pairs on that line would have to meet in other directions. Each other
direction can repair at most one such pair through its lone merged
block, since each of its two merged lines meets the split line once.
There are only five other directions. Thus five-point traces
cannot split. A four-point trace can only split into two pairs.

If no points were deleted, every direction has at most five blocks.
If one point was deleted, each direction has one
four-point trace and needs to split it without merging to obtain six
blocks. With no merges in any direction, none of these splits can
preserve intersection. These cases exhaust the possibilities.
\end{proof}

\section{Verification}

The independent replay generates degree profiles using combinations
with replacement, and obtains row bounds by explicitly enumerating
bounded integer tuples. A second program generates partitions
recursively and uses convex marginal costs for row bounds, with a
min-cost-flow cross-check on the surviving profile systems.
Both give five systems before the degree-three cut, three before
the congruence, and zero after it. The 13-edge residual and the
intersection checks use integer arithmetic. For 23 edges, the replay
regenerates the profile systems, checks each solver-proposed pair bound
as a rational inequality, checks the saved row certificates, and
enumerates and verifies the three final mask certificates.
For 24 edges we independently enumerate the residual degree profiles
in the degree-ten obstruction. The rotation and affine extension
checks use integer bit masks and exhaustive branching.
The earlier 22-edge lower bound appears in \cite{hrk}.

\begin{thebibliography}{9}
\bibitem{abw}
R. Aharoni, J. Bar\'at, and I. M. Wanless.
\emph{Multipartite hypergraphs achieving equality in Ryser's conjecture}.
arXiv:1409.4833v2, 2015.
\bibitem{akp}
A. Abu-Khazneh and A. Pokrovskiy.
\emph{Intersecting extremal constructions in Ryser's conjecture
for $r$-partite hypergraphs}. arXiv:1409.4938, 2014.
\bibitem{fhmw}
N. Franceti\'c, S. Herke, B. D. McKay, and I. M. Wanless.
\emph{On Ryser's conjecture for linear intersecting multipartite
hypergraphs}. arXiv:1508.00951, 2015.
\bibitem{hrk}
HRK44. \emph{The 22-edge lower bound for rank-six counterexamples
to Ryser's conjecture}. GitHub repository \texttt{HRK44/ryser-conjecture-22}, 2026.
\end{thebibliography}
\end{document}
